





\subsection{Related Work} \label{sec:related}
%\textbf{Related Work.}
%Safety has long been of interest in RL research. %Unlike traditional RL, 
Safe RL aims to learn policies that maximize expected reward on a task while respecting safety constraints during both learning and deployment \cite{garcia2015comprehensive}. 
Existing methods can be roughly classified as \emph{objective-based} or \emph{exploration-based}, depending on how safety is formulated.
We first discuss these categories, then the specific case of mobile robot navigation, which we use to evaluate our proposed method.

Objective-based methods encourage safety by penalizing constraint violations in the objective.
This can be done by relating cumulative reward to the system's risk, such as the probability of visiting error states \cite{geibel2005risk}.
In practice, this results in an RL agent attempting to minimize an empirical risk measure (that is, an approximation of the probability of entering a dangerous or undesired state).
Similarly, one can penalize the probability of losing reward (by visiting an unsafe state) for a given action \cite{mihatsch2002risk}, in which case the agent minimizes temporal differences in the reward and thus also minimizes risk.
Another approach is to restrict policies to be ergodic with high probability, meaning any state can eventually be reached from any other state \cite{moldovan2012safe}.
This is a more general problem, which comes at a cost: feasible safe policies do not always exist, and the algorithms are far more complex. 
While these methods can make an agent prefer safe actions, they cannot guarantee safety during training or deployment.
Another group of objective-based algorithms aims to modify the Markov Decision Process (MDP) that the RL agent tries to optimize.
Some model safe optimization problems as maximizing an unknown expected reward function \cite{pmlr-v37-sui15}.
However, they exploit regularity assumptions on the function wherein similar decisions are associated with similar rewards.
They also assume the bandit setting, where decisions do not cause state transitions.
% Others utilize constrained MDPs to enforce safety \cite{altman1999constrained}.
% These can be broadly classified as either offline or online approaches.
% Offline approaches \cite{achiam2017constrained,bharadhwaj2021conservative} collect data first, then perform the optimization.
% These methods can be quite conservative and hard to scale.
% Online methods \cite{BORKAR2005207,chow2017riskconstrained,bohez2019value}, on the other hand, couple both the data collection and optimization.
% However, they do not guarantee safety in training.
Others utilize constrained MDPs \cite{altman1999constrained} to enforce safety in various RL settings, either online or offline.
Online methods learn by coupling the iteration of numerical optimization algorithms (such as primal-dual gradient updates) with data collection \cite{BORKAR2005207,chow2017riskconstrained,tessler2018reward,bohez2019value}.
These algorithms have also been studied in exploration-based settings \cite{ding2020provably,efroni2020explorationexploitation}.
However, they provide no guarantees on safety during training.
On the other hand, offline schemes separate optimization and data collection \cite{achiam2017constrained,bharadhwaj2021conservative}.
They conservatively enforce safety constraints on every policy iteration but are more challenging to scale up.


Exploration-based methods modify the agent's exploration process instead of its optimization criterion.
Exploration is the process of learning about unexplored states by trying random actions or actions that are not expected to yield maximum reward (for example, an $\epsilon$-greedy strategy).
However, visiting unexplored states na\"ively can harm a robot or its environment.
To avoid this, one can aim to guarantee safety during both exploration and exploitation, in both training and testing, by modifying the exploration strategy to incorporate risk metrics \cite{riskmetricsRL}.
One can also use prior knowledge as an inductive bias for the exploration process \cite{garcia2015comprehensive,neuroevolutionaryinproceedings}; for example, one can provide a finite set of demonstrations as guidance on the task \cite{abbeel2005exploration}.
\tb{
Other approaches use control theory to guide or constrain the RL agent's actions.
Most of these approaches use system models with Lyapunov or control barrier functions (CBFs) to guarantee safety (or stability) of the system \cite{molnar2021model, liu2022robot, chow2018lyapunov}.
One can also combine data-driven approaches with model-free barrier or intervention functions \cite{cheng2019end, modelfreecbfsIVM, wagener2021safe}, or use robust CBFs to model uncertain parts of a system \cite{robustCBFarticle, cbfforunmodeleddynamics}.
Although these approaches can provide strong guarantees, most assume the system is control-affine, and need prior knowledge on some or all of the system model \cite{saferlforchanginglanes,conf:modelBasedReachRL,LewEtAl2021}, which may not always be feasible.
Finally, our method is most similar to \cite{dalal2018safe}, which learns a safety signal from data offline, then uses it to adjust an RL agent's controls at runtime.
This method uses a first-order safety approximation for fast adjustment, but (as we show) can be conservative.
}
 
%\subsection{Robot Navigation using Safe RL }
\tb{
Safe navigation is a fundamental challenge in robotics, because robots typically have uncertain, nonlinear dynamics.
Classical techniques such as A$^*$ or RRT \cite[Ch. 5]{lavalle2006planning} have been proposed to solve the navigation problem without learning.
With these methods, safety has been enforced at different levels of the planning hierarchy, such as trajectory planning \cite{kousik2020bridging}, or low-level control \cite{leung2020infusing}.
More recently, however, learning-based methods have been proposed \cite{saferlforchanginglanes,conf:modelBasedReachRL,LewEtAl2021}.
Some safe RL navigation approaches depend on learning a value function of the expected time that an agent takes to reach the desired goal \cite{chen2016decentralized, chen2018socially}.
Other approaches depend on learning the actions of the robot in an end-to-end manner \cite{long2018optimally, e2efirst, tai2018socially}, meaning that the agent attempts to convert raw sensor inputs (e.g., camera or LIDAR) into actuator commands.
The key advantage of RL over traditional planners is in accelerating computation time and solution quality. %by learning from past experience.
}